\documentclass[11pt]{article}
\usepackage{amsmath,amssymb}
\usepackage[margin=1in]{geometry}
\usepackage{hyperref}
\emergencystretch=3em

\title{A formal counterexample: mixed ratios evaluate on the face, not at the corner}
\author{Claude, with Daniel Murfet}
\date{2026-09-07}

\begin{document}
\maketitle
\sloppy

\begin{abstract}
Unit 206 records, as theorems, the smallest instance of the paper-facing correction reported since
unit~186: for \emph{mixed} exponent ratios the leading coefficient of a dressed normal moment
integral is the face functional of the amplitude, not its value at the origin. With $d=2$,
$h=(0,2)$, $k=(1,1)$, $\beta=1$ (ratios $\tfrac12,\tfrac32$; $\lambda=\tfrac12$, $J=\{0\}$) and
$\eta(u)=1+u_1$,
\[
\frac{\int_{(0,1]^2}(1+u_1)\,u_1^2\,e^{-Nu_0^2u_1^2}\,du}{N^{-1/2}}\longrightarrow\frac{5\sqrt\pi}{12},
\qquad
\eta(0)\cdot\frac{\Gamma(\tfrac12)}{2\,(2+1-1)}=\frac{\sqrt\pi}{4},
\]
and $5\sqrt\pi/12\ne\sqrt\pi/4$ (\texttt{MixedCounterexample.tendsto},
\texttt{corner\_evaluation\_fails}). Review v10 (units 204--205) reported no statement-level
defects; Astra~\#22 ranks the remaining work B2 (finite chart assembly) $\to$ A2 (general-$d$
stochastic transfer by equi-Lipschitz control) $\to$ D2 (signed orthants with phase, where the
phase reflects with $\varepsilon_\sigma=\prod_i\sigma_i^{k_i}$) $\to$ hand-off; C2 stays NO-GO.
Zero \texttt{sorry}/\texttt{axiom}; 9083 jobs.
\end{abstract}

\section{The things formalised}
{\small
\begin{verbatim}
namespace MixedCounterexample
def hEx : Fin 2 → ℕ := ![0, 2];  def kEx : Fin 2 → ℕ := ![1, 1]
noncomputable def ηEx (u : Fin 2 → ℝ) : ℝ := 1 + u 1
theorem ratio0 : ratioExp hEx kEx 0 = 1/2;  theorem ratio1 : ratioExp hEx kEx 1 = 3/2
theorem mixedConst_eq : monomialMixedConst hEx kEx (1/2) 1 = √π / 4
theorem mixedConst_shift_eq :
    monomialMixedConst (fun i => hEx i + ![0, 1] i) kEx (1/2) 1 = √π / 6
theorem amplitudeCoeff_eq : amplitudeCoeff hEx kEx (1/2) 1 ηEx = 5 * √π / 12
theorem corner_eq : ηEx 0 * monomialMixedConst hEx kEx (1/2) 1 = √π / 4
theorem corner_evaluation_fails :
    amplitudeCoeff hEx kEx (1/2) 1 ηEx ≠ ηEx 0 * monomialMixedConst hEx kEx (1/2) 1
theorem tendsto : Tendsto (fun N => (∫ x in unitBox 2, ηEx x * ((∏ i, x i ^ hEx i)
    * exp (-(1 * N * ∏ i, x i ^ (2 * kEx i))))) / N ^ (-(1/2 : ℝ))) atTop (𝓝 (5 * √π / 12))
\end{verbatim}
}

\section{Mathematics}
No integration is performed by hand: the amplitude $1+u_1$ is the sum of the two monomials
$u^{(0,0)}$ and $u^{(0,1)}$, and the face-moment theorem of unit~186
(\texttt{face\_moment\_eq}: $\int(\pi u)^\gamma\,\mathrm{faceLeadConst}\,w
=\mathrm{monomialMixedConst}(h+\gamma)$ when $\gamma$ vanishes on $J$) evaluates each as a bare
mixed constant, $\sqrt\pi/4$ for $h=(0,2)$ and $\sqrt\pi/6$ for $h=(0,3)$, via the two-dimensional
closed form of unit~198. The sum $5\sqrt\pi/12$ is the face value
$\tfrac{\sqrt\pi}2\int_0^1(1+v)v\,dv$; the corner value $\sqrt\pi/4$ forgets the $u_1$-integration
against the residual weight $u_1^{h_1-2k_1\lambda}=u_1$.

\section{Lean notes}
\begin{itemize}
\item \texttt{simp [ratioExp, hEx, kEx]} closes $(0+1)/(2\cdot1)=1/2$ outright but leaves
  $2+1=3$ for $(2+1)/(2\cdot1)=3/2$; a trailing \texttt{norm\_num} is needed in one case and is an
  error (``no goals'') in the other. Check each.
\item A hypothesis stated for the abbreviations \texttt{hEx kEx} is accepted by defeq as a term
  argument but not by \texttt{rw} inside a goal stated with \texttt{![0,2] ![1,1]}; use
  \texttt{show ratioExp hEx kEx 1 ≠ 1/2} first.
\item \texttt{Real.Gamma\_one\_half\_eq : Γ(1/2) = √π} and \texttt{Real.one\_rpow} reduce the
  closed forms; \texttt{nlinarith} with \texttt{0 < √π} separates $5\sqrt\pi/12$ from $\sqrt\pi/4$.
\end{itemize}

\section{Review v10 and Astra \#22}
Review v10 (units 204--205): clean; mirror should say ``$c>0$ on the closed cube'' with continuity,
keep the ``leading quotient limit, not full expansion'' caveat, and make clear that the cutoff face
integral is the unit-box parametrisation (the $b$-factor is displayed once). Astra~\#22: next is
B2, the conditional finite chart assembly (minimum exponent wins, then \emph{maximum} log
multiplicity; assemble numerator and denominator separately; 2--4 units), then A2 (general-$d$
stochastic transfer via eventual equi-Lipschitz dependence of the normalised integral on
$(\eta,\xi)$ in sup norm, 4--8 units), then D2 if the paper integrates full signed fibres (the phase
reflects: $\xi_\sigma(u)=\varepsilon_\sigma\xi(\sigma u)$, $\varepsilon_\sigma=\prod\sigma_i^{k_i}$;
odd monomials give $J_p(a)+J_p(-a)$, not $2J_p(a)$).
\end{document}
